\section{Introduction}
\label{sec:introduction}

The economic value of the equity of a bank, and the net interest income it earns, depend on the schedule on which every asset and every liability pays or reprices. For the contractual items that schedule is written in the contract. For the largest liability of a deposit-funded bank, the balances that customers may draw down at any time and that carry a rate the bank sets, it is not: the bank has to assume how much of the balance is stable, for how long, and how fast the deposit rate follows the market rate. The Basel standard for interest rate risk in the banking book lets a bank treat part of such a balance as a stable core with a long maturity, within caps by type of depositor \citep{bcbs2016,basel2019srp31}, and the banks that use their own behavioural models fit these assumptions on their own history, for example as a replicating portfolio of fixed-rate tranches whose weights minimise the variance of the margin between the portfolio yield and the client rate, a technique that the standard describes \citep{bcbs2016}; arbitrage-free and risk-management models of the deposit are the alternative \citep{jarrow1998,kalkbrener2004}. The assumption is consequential because a longer assumed maturity for a liability reduces the measured loss of economic value when rates rise, which is the direction in which a bank that holds long-duration assets is exposed. It is also a choice that the bank can change, and a limit on the loss of economic value is what makes the change attractive at the moment the limit is exceeded.

The review of the Federal Reserve into the failure of Silicon Valley Bank states exactly this practice as a finding. The bank had breached its long-term interest rate risk limits on and off since 2017, because of the structural mismatch between long-duration securities and short-duration deposits, and in April 2022 it made counterintuitive modelling assumptions about the duration of deposits to address the limit breach instead of managing the actual risk \citep{fed2023}. Related studies of the same period document the size of the exposure that the assumptions were hiding, the mark-to-market losses of US banks after the rise in rates \citep{jiang2024}, and the reason that deposit duration is a difficult quantity to assume, since the rates paid on deposits adjust slowly to the market rate and deposits flow out when it rises \citep{drechsler2017,drechsler2021}. A supervisor can read a finding of this kind after the failure. An independent validator of the model has to apply the check before it, on the history of the parameters that the bank has used, and for that a check has to be a procedure with an input, a decision rule and a known behaviour.

We found no such procedure in the standards and guidance on which the check would rest, nor in the literature on the models of non-maturing deposits that we cite below; the supervisory review states the practice as a finding and no document that we found turns it into a test. That absence is a statement about the documents we read, and we do not claim more. This work proposes the procedure and measures what it does. The proposed test takes the history of one behavioural parameter, the dates, and a flag at each date that says whether a limit was breached in the window ending there, and counts the dates at which the parameter changed while the flag was set. It returns red when the count is at least one, green when it is zero, and no number when there are fewer than two dates. It has no threshold to tune and it is implemented in the engine of the CAOS\_Contraste product as a function of the validation module for interest rate risk \citep{riskvalidation2026}. The product applies it to the wholesale deposits of the bank of its case C14, whose call-report aggregates are real and whose behavioural assumptions are synthetic.

The test is deliberately simple, and the questions that matter are what a count of coincidences can and cannot establish. We answer them with the committed results of the product, with constructed and exhaustively enumerated histories, with simulations that plant an assumption change at breach dates in two kinds of parameter, and with proofs. The two kinds are the parameters that the bank states and holds until it decides to change them, and the parameters that are re-estimated at every date from the data, such as the weighted average maturity of a replicating portfolio. The difference between them turns out to decide what the test can say.

\paragraph{Contributions.}
\begin{enumerate}
\item A precise statement of the test as a statistic, a decision rule and a family with a tolerance for the size of a change, with the algorithm and the implementation in the engine of the product (Section~\ref{sec:proposed-method}).
\item Eight propositions with proofs (Section~\ref{sec:theory}): the test depends on the parameter only through the indicator of a change; on a stated parameter it flags every change made at a breach date and nothing else; its false-flag probability under unrelated changes is $1-(1-\lambda)^m$ or, for uniformly placed dates, a hypergeometric expression; on a parameter whose changes have no atom it is red whenever a breach exists, whatever the parameter did; the one-sided Fisher test is its permutation test; and a tolerance calibrated by a rank rule holds its size in finite samples.
\item The reproduction of the committed rows of the product by the engine, the enumeration of every history on a binary grid and of 
andomN{} random histories against the closed-form count, and the demonstration that the last digit of a floating-point value is a change (Section~\ref{sec:results}).
\item A measurement, on \statedN{} histories per arm in each of \statedCells{} cells, of the false-flag and detection rates of the test and of four published tests run on the same histories, and a measurement on \fitN{} simulated banks and on the real fitted series of the bank for a re-estimated parameter (Section~\ref{sec:results}).
\item A negative result that bounds the claim: on a re-estimated parameter the test flags \fitNullEngine{} of null histories and its area under the ROC curve against a planted 2-year lengthening is \fitAucCEngine{}, so a red result there carries no information about the bank; a tolerance rule calibrated on a null population restores the discrimination at the sizes of override that exceed the natural changes (Sections~\ref{sec:results} and~\ref{sec:ablation}).
\end{enumerate}

\paragraph{Main result.}
On a stated parameter the proposed test flagged all of the \statedN{} histories with a planted change at a breach date in the base cell (lower bound \statedAltLoEngine{}\,\%), and it also flagged \statedNullEngine{} of the histories whose unrelated changes happened to coincide with a breach, a rate that matches the closed form of Proposition~\ref{prop:false-flag} in \bernInside{} of \bernCells{} cells. On a re-estimated parameter it flagged \fitNullEngine{} of the null banks and \fitAddCEngine{} of those with a 2-year override. The tolerance calibrated at 5\,\% on a separate population flagged \fitNullTol{} of the null banks and \fitAddDTol{} of the banks with a 4-year override, and \fitAddCTol{} of those with a 2-year override. Of the nine predictions we tested, \hypothesesHeld{} held and \hypothesesRefuted{} were refuted (Table~\ref{tab:hypotheses}).

The rest of the paper is organised as follows. Section~\ref{sec:related-work} places the test among the supervisory documents, the models of deposits, the studies of the manipulation of regulatory measures and the methods of change detection. Section~\ref{sec:problem} formulates the problem and Section~\ref{sec:proposed-method} states the method. Section~\ref{sec:theory} proves its properties. Sections~\ref{sec:experimental-design} to~\ref{sec:ablation} describe and report the experiments and the ablation. Sections~\ref{sec:discussion} to~\ref{sec:conclusions} discuss the results, their limits, the reproduction of every number and the conclusions.
